\chapter{Θεωρητικό υπόβαθρο}

% TODO(PR13): Πλήρης συγγραφή, με βάση τα κεφάλαια 2, 3, 6 και 10 του
% \textcite{han2011data}.

\section{Προεπεξεργασία δεδομένων}

% Κεφάλαια 2 και 3 του βιβλίου: καθαρισμός, ολοκλήρωση, μετασχηματισμός και
% μείωση δεδομένων.

\section{Εξόρυξη συχνών υποσυνόλων}

% Κεφάλαιο 6: ορισμοί υποστήριξης και εμπιστοσύνης, αλγόριθμος Apriori,
% αλγόριθμος FP-Growth, κανόνες συσχέτισης και το μέτρο lift.

\section{Συσταδοποίηση}

% Κεφάλαιο 10: k-means, ιεραρχική συσταδοποίηση, μέτρα αξιολόγησης.

\section{Ανάλυση RFM και τμηματοποίηση πελατών}

% Το μοντέλο RFM και η εφαρμογή του σε δεδομένα λιανικής, με αναφορά στους
% \textcite{chen2012data}.

\section{Οικονομικοί δείκτες ηλεκτρονικού εμπορίου}

% GMV, ποσοστό επαναληπτικών αγορών, διατήρηση πελατών ανά κοόρτη, CLV, CAC
% και ο λόγος CLV:CAC.
